\documentclass[11pt,a4paper]{article}
\usepackage[margin=18mm,headheight=16pt]{geometry}
\usepackage{fontspec}
\setmainfont{DejaVu Sans}
\usepackage{microtype,parskip,setspace,ragged2e}
\usepackage{xcolor,graphicx,enumitem,fancyhdr}
\usepackage[most]{tcolorbox}
\usepackage[colorlinks=true,linkcolor=teal!60!black,urlcolor=teal!60!black]{hyperref}
\definecolor{teal}{HTML}{187C80}
\definecolor{pale}{HTML}{EEF6F5}
\definecolor{ink}{HTML}{203342}
\definecolor{accent}{HTML}{A62C40}
\setstretch{1.12}
\setlength{\parindent}{0pt}
\setlist[itemize]{leftmargin=1.3em,itemsep=0.25em,topsep=0.2em,parsep=0.1em}
\pagestyle{fancy}
\fancyhf{}
\lhead{\small ML Architecture --- Speaker Notes}
\rhead{\small Bank customer-churn example}
\cfoot{\thepage}
\renewcommand{\headrulewidth}{0.3pt}
\tcbset{breakable,arc=1.5mm,boxrule=0.6pt,left=3mm,right=3mm,top=2mm,bottom=2mm,before skip=6pt,after skip=6pt,fonttitle=\bfseries\small,before upper={\setlength{\parskip}{0.4em}\RaggedRight}}
\newtcolorbox{saybox}{colback=pale,colframe=teal,title={TO SAY},fontupper=\fontsize{11}{14.5}\selectfont}
\newtcolorbox{cuebox}[1][]{colback=accent!5,colframe=accent!70!black,title={CUES},#1}
\begin{document}
\thispagestyle{empty}
\begin{center}
{\Huge\bfseries ML Architecture}\\[4mm]
{\Large Speaker Notes}\\[2mm]
{\large A bank customer-churn case --- 20 October 2026}\\[8mm]
\end{center}
\noindent Each page shows one slide exactly as presenters will see it, the words to speak, and short delivery cues. The script is a speaking aid, not a transcript to read verbatim --- adapt it to your own voice. Animated builds are marked; press the arrow key the stated number of times before moving to the next slide.
\clearpage

\section*{\normalsize 01 / 55 \quad ML Architecture}
\noindent\textcolor{teal!70!black}{\small Block 1 / Foundations \textbullet\ All presenters}\par\vspace{2mm}
\noindent\begin{minipage}[t]{0.40\textwidth}\vspace{0pt}
\fcolorbox{teal!30}{white}{\includegraphics[width=\dimexpr\linewidth-2\fboxsep-2\fboxrule\relax]{previews/speaker-notes/slide-01.png}}
\end{minipage}\hfill
\begin{minipage}[t]{0.57\textwidth}\vspace{0pt}
\begin{cuebox}
\begin{itemize}
\item New block: \textbf{Block 1 / Foundations}.
\item Presenter: \textbf{All presenters}.
\item Pause briefly before the bank example — let the concept land first.
\item Close on the takeaway line before advancing.
\end{itemize}
\end{cuebox}
\end{minipage}
\vspace{2mm}
\begin{saybox}
From a trained model to a maintainable, observable, governed system. \textbf{Software engineering} for systems that depend on learned behavior. \textbf{One shared case}: bank customer-churn prediction. \textbf{Two blocks}: model-to-system foundations, then architecture and quality. \textbf{Five presenters}: Ricardo, Tomás, Simon, Luca, and Matheus. Bank churn supplies the shared case. Examples, records, schedules, and incidents are illustrative. \textit{A useful model needs an engineered system around it.}
\end{saybox}
\par\noindent{\scriptsize\textcolor{teal!60!black}{Reading anchor: Bass, Ch. 14 · Kalinowski, Ch. 13–14}}
\clearpage
\section*{\normalsize 02 / 55 \quad Two blocks, one shared bank case}
\noindent\textcolor{teal!70!black}{\small Block 1 / Foundations \textbullet\ All presenters}\par\vspace{2mm}
\noindent\begin{minipage}[t]{0.40\textwidth}\vspace{0pt}
\fcolorbox{teal!30}{white}{\includegraphics[width=\dimexpr\linewidth-2\fboxsep-2\fboxrule\relax]{previews/speaker-notes/slide-02.png}}
\end{minipage}\hfill
\begin{minipage}[t]{0.57\textwidth}\vspace{0pt}
\begin{cuebox}
\begin{itemize}
\item Presenter: \textbf{All presenters}.
\item Pause briefly before the bank example — let the concept land first.
\item Close on the takeaway line before advancing.
\end{itemize}
\end{cuebox}
\end{minipage}
\vspace{2mm}
\begin{saybox}
Follow the group’s topic outline; use the case to connect each architectural decision. \textbf{Block 1}: lifecycle, serving, monitoring, and design decisions. \textbf{Block 2}: tactics, MLOps, security, privacy, and quality. \textbf{Concept first}: explain the mechanism, then apply it to churn. \textbf{Named handoffs}: each section identifies its presenter; unassigned topics remain TBD. \textit{At our bank,} the same source → feature → model → score → policy path returns throughout the lecture. \textit{Each architectural decision connects predictive behavior to a system responsibility.}
\end{saybox}
\par\noindent{\scriptsize\textcolor{teal!60!black}{Reading anchor: Group discussion document · Agreed lecture redesign}}
\clearpage
\section*{\normalsize 03 / 55 \quad The bank case: features, labels, and predictions}
\noindent\textcolor{teal!70!black}{\small Block 1 / Foundations \textbullet\ Ricardo}\par\vspace{2mm}
\noindent\begin{minipage}[t]{0.40\textwidth}\vspace{0pt}
\fcolorbox{teal!30}{white}{\includegraphics[width=\dimexpr\linewidth-2\fboxsep-2\fboxrule\relax]{previews/speaker-notes/slide-03.png}}
\end{minipage}\hfill
\begin{minipage}[t]{0.57\textwidth}\vspace{0pt}
\begin{cuebox}
\begin{itemize}
\item Presenter: \textbf{Ricardo}.
\item Animated build: click 2 times to reveal the remaining steps before moving on.
\item Pause briefly before the bank example — let the concept land first.
\item Close on the takeaway line before advancing.
\end{itemize}
\end{cuebox}
\end{minipage}
\vspace{2mm}
\begin{saybox}
A few ML terms connect familiar application data to a learned risk estimate. \textbf{Feature}: a prepared input, such as logins during the previous seven days. \textbf{Label}: the observed churn outcome used for training and later evaluation. \textbf{Prediction}: an estimate produced before the customer’s future outcome is known. \textbf{Training vs. inference}: learn from earlier examples; apply the approved model to current inputs. \textit{At our bank,} ana / C042: logins fall from 7 to 1; the illustrative score is 0.70. Her next-30-day outcome is unknown. \textit{Lower activity is evidence for a model input, not proof that a customer will leave.}
\end{saybox}
\par\noindent{\scriptsize\textcolor{teal!60!black}{Reading anchor: Bass, Ch. 14 · Simon’s notes, §1}}
\clearpage
\section*{\normalsize 04 / 55 \quad From a trained model to an ML system}
\noindent\textcolor{teal!70!black}{\small Block 1 / Foundations \textbullet\ Ricardo}\par\vspace{2mm}
\noindent\begin{minipage}[t]{0.40\textwidth}\vspace{0pt}
\fcolorbox{teal!30}{white}{\includegraphics[width=\dimexpr\linewidth-2\fboxsep-2\fboxrule\relax]{previews/speaker-notes/slide-04.png}}
\end{minipage}\hfill
\begin{minipage}[t]{0.57\textwidth}\vspace{0pt}
\begin{cuebox}
\begin{itemize}
\item Presenter: \textbf{Ricardo}.
\item Animated build: click 4 times to reveal the remaining steps before moving on.
\item Pause briefly before the bank example — let the concept land first.
\item Close on the takeaway line before advancing.
\end{itemize}
\end{cuebox}
\end{minipage}
\vspace{2mm}
\begin{saybox}
Architecture connects learning, production data, operational controls, and business decisions. \textbf{Training} uses historical features and labels to build and evaluate candidates. \textbf{Production} validates current data and applies an approved, compatible model. \textbf{Stored scores} link predictions to customer, time, input evidence, and versions. \textbf{Business policy} determines the permitted action; it remains separate from the score. \textit{At our bank,} the bank computes a monthly batch of churn scores, then applies campaign eligibility and contact rules. \textit{Keep the learning path and production path distinct, but traceable.}
\end{saybox}
\par\noindent{\scriptsize\textcolor{teal!60!black}{Reading anchor: Bass, Ch. 14 · Kalinowski, §13.2–13.3}}
\clearpage
\section*{\normalsize 05 / 55 \quad Two lifecycles, one compatible system}
\noindent\textcolor{teal!70!black}{\small Block 1 / Foundations \textbullet\ Ricardo}\par\vspace{2mm}
\noindent\begin{minipage}[t]{0.40\textwidth}\vspace{0pt}
\fcolorbox{teal!30}{white}{\includegraphics[width=\dimexpr\linewidth-2\fboxsep-2\fboxrule\relax]{previews/speaker-notes/slide-05.png}}
\end{minipage}\hfill
\begin{minipage}[t]{0.57\textwidth}\vspace{0pt}
\begin{cuebox}
\begin{itemize}
\item Presenter: \textbf{Ricardo}.
\item Animated build: click 3 times to reveal the remaining steps before moving on.
\item Pause briefly before the bank example — let the concept land first.
\item Close on the takeaway line before advancing.
\end{itemize}
\end{cuebox}
\end{minipage}
\vspace{2mm}
\begin{saybox}
Code correctness and predictive usefulness depend on different feedback and release cycles. \textbf{Software changes} can be checked through code, contract, and integration tests. \textbf{Model development} also needs representative data, suitable labels, and evaluation. \textbf{Feedback differs}: execution checks may be immediate; churn labels arrive later. \textbf{Co-versioning} links changes across code, feature meaning, model, and configuration. \textit{At our bank,} a CRM connector can change today while M1 stays fixed; its changed inputs can still alter predictions. \textit{A passing software test does not establish future predictive quality.}
\end{saybox}
\par\noindent{\scriptsize\textcolor{teal!60!black}{Reading anchor: Kalinowski, Ch. 13–14 · Bass, Ch. 11 · Simon’s notes, §1}}
\clearpage
\section*{\normalsize 06 / 55 \quad Prediction and business action are separate}
\noindent\textcolor{teal!70!black}{\small Block 1 / Foundations \textbullet\ Ricardo}\par\vspace{2mm}
\noindent\begin{minipage}[t]{0.40\textwidth}\vspace{0pt}
\fcolorbox{teal!30}{white}{\includegraphics[width=\dimexpr\linewidth-2\fboxsep-2\fboxrule\relax]{previews/speaker-notes/slide-06.png}}
\end{minipage}\hfill
\begin{minipage}[t]{0.57\textwidth}\vspace{0pt}
\begin{cuebox}
\begin{itemize}
\item Presenter: \textbf{Ricardo}.
\item Animated build: click 3 times to reveal the remaining steps before moving on.
\item Pause briefly before the bank example — let the concept land first.
\item Close on the takeaway line before advancing.
\end{itemize}
\end{cuebox}
\end{minipage}
\vspace{2mm}
\begin{saybox}
A model estimates risk. Business rules decide whether and how the bank may act. \textbf{Score}: estimated churn risk, not a guaranteed outcome. \textbf{Policy}: preferences, eligibility, contact limits, and campaign budget. \textbf{Action}: an offer, another response, or no contact at all. \textbf{Separate evaluation}: predictive quality and the benefit of an offer are different questions. \textit{At our bank,} ana may have a high risk score and still receive no offer because contact preferences or policy prohibit it. \textit{A policy change should not require retraining the prediction model.}
\end{saybox}
\par\noindent{\scriptsize\textcolor{teal!60!black}{Reading anchor: Bass, Ch. 14 · Teaching policy decomposition}}
\clearpage
\section*{\normalsize 07 / 55 \quad Serving patterns: when is a prediction needed?}
\noindent\textcolor{teal!70!black}{\small Block 1 / Deployment \& serving \textbullet\ Tomás}\par\vspace{2mm}
\noindent\begin{minipage}[t]{0.40\textwidth}\vspace{0pt}
\fcolorbox{teal!30}{white}{\includegraphics[width=\dimexpr\linewidth-2\fboxsep-2\fboxrule\relax]{previews/speaker-notes/slide-07.png}}
\end{minipage}\hfill
\begin{minipage}[t]{0.57\textwidth}\vspace{0pt}
\begin{cuebox}
\begin{itemize}
\item New block: \textbf{Block 1 / Deployment \& serving}.
\item Presenter: \textbf{Tomás}.
\item Animated build: click 2 times to reveal the remaining steps before moving on.
\item Pause briefly before the bank example — let the concept land first.
\item Close on the takeaway line before advancing.
\end{itemize}
\end{cuebox}
\end{minipage}
\vspace{2mm}
\begin{saybox}
Select the interaction pattern from the decision’s timing and failure behavior. \textbf{Batch}: compute many predictions; consumers read stored results later. \textbf{Synchronous}: the caller waits for a result during its request. \textbf{Asynchronous}: an event producer and consumer are decoupled through messaging. \textbf{Edge}: a deployment location, independent of prediction timing. \textit{At our bank,} the churn campaign can use precomputed scores. An immediate fraud check motivates a different timing need. \textit{Timing, location, and release strategy are separate choices.}
\end{saybox}
\par\noindent{\scriptsize\textcolor{teal!60!black}{Reading anchor: Kalinowski, §13.3–13.4; supporting Ch. 12}}
\clearpage
\section*{\normalsize 08 / 55 \quad Batch serving: validate, score, publish}
\noindent\textcolor{teal!70!black}{\small Block 1 / Deployment \& serving \textbullet\ Tomás}\par\vspace{2mm}
\noindent\begin{minipage}[t]{0.40\textwidth}\vspace{0pt}
\fcolorbox{teal!30}{white}{\includegraphics[width=\dimexpr\linewidth-2\fboxsep-2\fboxrule\relax]{previews/speaker-notes/slide-08.png}}
\end{minipage}\hfill
\begin{minipage}[t]{0.57\textwidth}\vspace{0pt}
\begin{cuebox}
\begin{itemize}
\item Presenter: \textbf{Tomás}.
\item Animated build: click 2 times to reveal the remaining steps before moving on.
\item Pause briefly before the bank example — let the concept land first.
\item Close on the takeaway line before advancing.
\end{itemize}
\end{cuebox}
\end{minipage}
\vspace{2mm}
\begin{saybox}
A complete validated batch must be available before its business deadline. \textbf{Snapshot}: define population, cutoff, and required source freshness. \textbf{Validate}: reject incompatible inputs before scoring. \textbf{Reconcile}: account for written, rejected, and withheld records. \textbf{Publish}: expose a validated batch under a defined consumer contract. \textit{At our bank,} monthly churn scores are saved for the campaign. A real numerical deadline must be agreed with campaign owners. \textit{Optimize the whole batch deadline, not only per-record prediction speed.}
\end{saybox}
\par\noindent{\scriptsize\textcolor{teal!60!black}{Reading anchor: Bass, Ch. 14 · Kalinowski, §13.3}}
\clearpage
\section*{\normalsize 09 / 55 \quad Synchronous serving: the caller waits}
\noindent\textcolor{teal!70!black}{\small Block 1 / Deployment \& serving \textbullet\ Tomás}\par\vspace{2mm}
\noindent\begin{minipage}[t]{0.40\textwidth}\vspace{0pt}
\fcolorbox{teal!30}{white}{\includegraphics[width=\dimexpr\linewidth-2\fboxsep-2\fboxrule\relax]{previews/speaker-notes/slide-09.png}}
\end{minipage}\hfill
\begin{minipage}[t]{0.57\textwidth}\vspace{0pt}
\begin{cuebox}
\begin{itemize}
\item Presenter: \textbf{Tomás}.
\item Animated build: click 2 times to reveal the remaining steps before moving on.
\item Pause briefly before the bank example — let the concept land first.
\item Close on the takeaway line before advancing.
\end{itemize}
\end{cuebox}
\end{minipage}
\vspace{2mm}
\begin{saybox}
Prediction becomes part of the application’s response path. \textbf{Response time} includes lookup, preparation, inference, and network work. \textbf{Availability} now affects the waiting application’s request. \textbf{Failure behavior} needs explicit timeout, error, and fallback policies. \textbf{Consistency} still requires compatible feature meanings and model versions. \textit{At our bank,} an immediate fraud check can justify this pattern. The monthly churn campaign does not require it. \textit{A fast model does not guarantee a fast end-to-end service.}
\end{saybox}
\par\noindent{\scriptsize\textcolor{teal!60!black}{Reading anchor: Kalinowski, §13.3–13.4 · Teaching comparison}}
\clearpage
\section*{\normalsize 10 / 55 \quad Asynchronous serving: decouple the work}
\noindent\textcolor{teal!70!black}{\small Block 1 / Deployment \& serving \textbullet\ Tomás}\par\vspace{2mm}
\noindent\begin{minipage}[t]{0.40\textwidth}\vspace{0pt}
\fcolorbox{teal!30}{white}{\includegraphics[width=\dimexpr\linewidth-2\fboxsep-2\fboxrule\relax]{previews/speaker-notes/slide-10.png}}
\end{minipage}\hfill
\begin{minipage}[t]{0.57\textwidth}\vspace{0pt}
\begin{cuebox}
\begin{itemize}
\item Presenter: \textbf{Tomás}.
\item Animated build: click 2 times to reveal the remaining steps before moving on.
\item Pause briefly before the bank example — let the concept land first.
\item Close on the takeaway line before advancing.
\end{itemize}
\end{cuebox}
\end{minipage}
\vspace{2mm}
\begin{saybox}
Messages separate production, processing, and consumption of a prediction. \textbf{Buffering} absorbs bursts and supports independent processing. \textbf{Lag} and queue health become operational concerns. \textbf{Retries} require failure handling and idempotent processing. \textbf{Contracts} define message identity, timing, schema, and result interpretation. \textit{At our bank,} if a future campaign consumes prediction events, a repeated message must not cause duplicate customer contacts. \textit{Decoupling adds messaging responsibilities; it does not remove timing requirements.}
\end{saybox}
\par\noindent{\scriptsize\textcolor{teal!60!black}{Reading anchor: Kalinowski, §13.3 · Bank application is hypothetical}}
\clearpage
\section*{\normalsize 11 / 55 \quad Deployment location is a separate decision}
\noindent\textcolor{teal!70!black}{\small Block 1 / Deployment \& serving \textbullet\ Tomás}\par\vspace{2mm}
\noindent\begin{minipage}[t]{0.40\textwidth}\vspace{0pt}
\fcolorbox{teal!30}{white}{\includegraphics[width=\dimexpr\linewidth-2\fboxsep-2\fboxrule\relax]{previews/speaker-notes/slide-11.png}}
\end{minipage}\hfill
\begin{minipage}[t]{0.57\textwidth}\vspace{0pt}
\begin{cuebox}
\begin{itemize}
\item Presenter: \textbf{Tomás}.
\item Pause briefly before the bank example — let the concept land first.
\item Close on the takeaway line before advancing.
\end{itemize}
\end{cuebox}
\end{minipage}
\vspace{2mm}
\begin{saybox}
Central, embedded, and edge deployments change operations and constraints. \textbf{Central}: simplify shared operation and governance. \textbf{Embedded}: deploy the model with its host application. \textbf{Edge}: account for constrained hardware and intermittent connectivity. \textbf{Choose independently}: location does not uniquely determine batch or request timing. \textit{At our bank,} keep churn centrally governed. Device-local inference is a contrast, not an extra subsystem to implement here. \textit{Architecture follows the operational context.}
\end{saybox}
\par\noindent{\scriptsize\textcolor{teal!60!black}{Reading anchor: Kalinowski, supporting Ch. 12; §13.3–13.4}}
\clearpage
\section*{\normalsize 12 / 55 \quad Package the whole compatible set}
\noindent\textcolor{teal!70!black}{\small Block 1 / Deployment \& serving \textbullet\ Tomás}\par\vspace{2mm}
\noindent\begin{minipage}[t]{0.40\textwidth}\vspace{0pt}
\fcolorbox{teal!30}{white}{\includegraphics[width=\dimexpr\linewidth-2\fboxsep-2\fboxrule\relax]{previews/speaker-notes/slide-12.png}}
\end{minipage}\hfill
\begin{minipage}[t]{0.57\textwidth}\vspace{0pt}
\begin{cuebox}
\begin{itemize}
\item Presenter: \textbf{Tomás}.
\item Animated build: click 2 times to reveal the remaining steps before moving on.
\item Pause briefly before the bank example — let the concept land first.
\item Close on the takeaway line before advancing.
\end{itemize}
\end{cuebox}
\end{minipage}
\vspace{2mm}
\begin{saybox}
A deployable release contains more than the learned model file. \textbf{Model artifact}: the learned parameters and required format. \textbf{Runtime}: dependencies, preparation code, and execution environment. \textbf{Contract}: feature meaning, types, missingness, and temporal cutoff. \textbf{Configuration}: scoring population and execution settings. \textit{At our bank,} the bank packages M1 with the exact feature interpretation needed to score the monthly customer snapshot. \textit{A package that loads successfully still needs release and quality evidence.}
\end{saybox}
\par\noindent{\scriptsize\textcolor{teal!60!black}{Reading anchor: Bass, Ch. 14 · Kalinowski, §13.3–13.4}}
\clearpage
\section*{\normalsize 13 / 55 \quad Shadow deployment: observe without acting}
\noindent\textcolor{teal!70!black}{\small Block 1 / Deployment \& serving \textbullet\ Tomás}\par\vspace{2mm}
\noindent\begin{minipage}[t]{0.40\textwidth}\vspace{0pt}
\fcolorbox{teal!30}{white}{\includegraphics[width=\dimexpr\linewidth-2\fboxsep-2\fboxrule\relax]{previews/speaker-notes/slide-13.png}}
\end{minipage}\hfill
\begin{minipage}[t]{0.57\textwidth}\vspace{0pt}
\begin{cuebox}
\begin{itemize}
\item Presenter: \textbf{Tomás}.
\item Animated build: click 4 times to reveal the remaining steps before moving on.
\item Pause briefly before the bank example — let the concept land first.
\item Close on the takeaway line before advancing.
\end{itemize}
\end{cuebox}
\end{minipage}
\vspace{2mm}
\begin{saybox}
A candidate scores real inputs while the approved model remains the action source. \textbf{Run in parallel} against comparable production inputs. \textbf{Record identity}: distinguish M1 and M2 outputs explicitly. \textbf{Compare} operational behavior and score differences first. \textbf{Wait for labels} to compare future churn quality. \textit{At our bank,} only the M1 record is eligible for campaign use; M2’s shadow score cannot create a second offer. \textit{Shadow checks observation, not the business effect of acting on the candidate.}
\end{saybox}
\par\noindent{\scriptsize\textcolor{teal!60!black}{Reading anchor: Group outline · Breck, infrastructure 6–7 · Teaching application}}
\clearpage
\section*{\normalsize 14 / 55 \quad Canary deployment: limit exposure}
\noindent\textcolor{teal!70!black}{\small Block 1 / Deployment \& serving \textbullet\ Tomás}\par\vspace{2mm}
\noindent\begin{minipage}[t]{0.40\textwidth}\vspace{0pt}
\fcolorbox{teal!30}{white}{\includegraphics[width=\dimexpr\linewidth-2\fboxsep-2\fboxrule\relax]{previews/speaker-notes/slide-14.png}}
\end{minipage}\hfill
\begin{minipage}[t]{0.57\textwidth}\vspace{0pt}
\begin{cuebox}
\begin{itemize}
\item Presenter: \textbf{Tomás}.
\item Animated build: click 3 times to reveal the remaining steps before moving on.
\item Pause briefly before the bank example — let the concept land first.
\item Close on the takeaway line before advancing.
\end{itemize}
\end{cuebox}
\end{minipage}
\vspace{2mm}
\begin{saybox}
An approved limited cohort uses the candidate before broader rollout. \textbf{Restrict exposure} to a defined cohort and scope. \textbf{Observe} execution, compatibility, coverage, and policy behavior. \textbf{Define stop rules} and a rehearsed recovery path. \textbf{Separate evidence}: short operational success cannot prove 30-day predictive quality. \textit{At our bank,} the bank may authorize a small cohort to use M2 while retaining M1 for the remainder. \textit{A canary controls risk; it is not automatically a valid business experiment.}
\end{saybox}
\par\noindent{\scriptsize\textcolor{teal!60!black}{Reading anchor: Breck, infrastructure test 6 · Bank teaching application}}
\clearpage
\section*{\normalsize 15 / 55 \quad A/B testing: define the intervention}
\noindent\textcolor{teal!70!black}{\small Block 1 / Deployment \& serving \textbullet\ Tomás}\par\vspace{2mm}
\noindent\begin{minipage}[t]{0.40\textwidth}\vspace{0pt}
\fcolorbox{teal!30}{white}{\includegraphics[width=\dimexpr\linewidth-2\fboxsep-2\fboxrule\relax]{previews/speaker-notes/slide-15.png}}
\end{minipage}\hfill
\begin{minipage}[t]{0.57\textwidth}\vspace{0pt}
\begin{cuebox}
\begin{itemize}
\item Presenter: \textbf{Tomás}.
\item Animated build: click 3 times to reveal the remaining steps before moving on.
\item Pause briefly before the bank example — let the concept land first.
\item Close on the takeaway line before advancing.
\end{itemize}
\end{cuebox}
\end{minipage}
\vspace{2mm}
\begin{saybox}
Assigned groups and a stated outcome make the comparison interpretable. \textbf{State what changes}: model, decision policy, or retention treatment. \textbf{Assign comparable groups} under an appropriate experiment design. \textbf{Define outcomes} and the observation window before interpreting results. \textbf{Avoid confounding}: changing model and offer together obscures their separate effects. \textit{At our bank,} compare two offer policies separately from asking whether M2 ranks churn risk better than M1. \textit{Prediction accuracy and the effect of contacting a customer are different targets.}
\end{saybox}
\par\noindent{\scriptsize\textcolor{teal!60!black}{Reading anchor: Group outline · Simon’s notes, §3 · Teaching application}}
\clearpage
\section*{\normalsize 16 / 55 \quad How the world changes over time}
\noindent\textcolor{teal!70!black}{\small Block 1 / Monitoring \& degradation \textbullet\ Simon}\par\vspace{2mm}
\noindent\begin{minipage}[t]{0.40\textwidth}\vspace{0pt}
\fcolorbox{teal!30}{white}{\includegraphics[width=\dimexpr\linewidth-2\fboxsep-2\fboxrule\relax]{previews/speaker-notes/slide-16.png}}
\end{minipage}\hfill
\begin{minipage}[t]{0.57\textwidth}\vspace{0pt}
\begin{cuebox}
\begin{itemize}
\item New block: \textbf{Block 1 / Monitoring \& degradation}.
\item Presenter: \textbf{Simon}.
\item Pause briefly before the bank example — let the concept land first.
\item Close on the takeaway line before advancing.
\item Adaptation note on file for this slide; do not read aloud unless asked.
\end{itemize}
\end{cuebox}
\end{minipage}
\vspace{2mm}
\begin{saybox}
Sudden, gradual, incremental, and recurring changes are different time patterns—not different input features. \textbf{Sudden}: the old regime switches abruptly. \textbf{Gradual}: old and new regimes coexist during a transition. \textbf{Incremental / recurring}: intermediate change differs from a returning regime. \textbf{Blips}: isolated noise is not necessarily a lasting regime change. Bank illustrations: a pricing change can be sudden; staggered adoption gradual; evolving habits incremental; seasonal activity recurring. \textit{First ask how change unfolds over time. Then ask what changed: input data, outcome relationships, or data validity.}
\end{saybox}
\par\noindent{\scriptsize\textcolor{teal!60!black}{Reading anchor: Krawczyk \& Cano (2018), Fig. 2 · Original figure; attribution and source link in guide}}
\clearpage
\section*{\normalsize 17 / 55 \quad Monitoring: observe more than model accuracy}
\noindent\textcolor{teal!70!black}{\small Block 1 / Monitoring \& degradation \textbullet\ Simon}\par\vspace{2mm}
\noindent\begin{minipage}[t]{0.40\textwidth}\vspace{0pt}
\fcolorbox{teal!30}{white}{\includegraphics[width=\dimexpr\linewidth-2\fboxsep-2\fboxrule\relax]{previews/speaker-notes/slide-17.png}}
\end{minipage}\hfill
\begin{minipage}[t]{0.57\textwidth}\vspace{0pt}
\begin{cuebox}
\begin{itemize}
\item Presenter: \textbf{Simon}.
\item Pause briefly before the bank example — let the concept land first.
\item Close on the takeaway line before advancing.
\end{itemize}
\end{cuebox}
\end{minipage}
\vspace{2mm}
\begin{saybox}
Inputs, system execution, predictions, and business actions need different evidence. \textbf{Inputs}: are the right records available and meaningful? \textbf{System}: did the pipeline complete correctly and on time? \textbf{Model}: are predictions useful on the current population? \textbf{Business}: are permitted actions and outcomes being tracked? \textit{At our bank,} a healthy churn job can still use a shifted population. A useful score can still lead to an inappropriate contact. \textit{Observability connects evidence to investigation; dashboards alone do not decide the remedy.}
\end{saybox}
\par\noindent{\scriptsize\textcolor{teal!60!black}{Reading anchor: Simon’s notes, §2–4 · Bass, Ch. 11 / §14.1.5}}
\clearpage
\section*{\normalsize 18 / 55 \quad Monitoring architecture: from signal to response}
\noindent\textcolor{teal!70!black}{\small Block 1 / Monitoring \& degradation \textbullet\ Simon}\par\vspace{2mm}
\noindent\begin{minipage}[t]{0.40\textwidth}\vspace{0pt}
\fcolorbox{teal!30}{white}{\includegraphics[width=\dimexpr\linewidth-2\fboxsep-2\fboxrule\relax]{previews/speaker-notes/slide-18.png}}
\end{minipage}\hfill
\begin{minipage}[t]{0.57\textwidth}\vspace{0pt}
\begin{cuebox}
\begin{itemize}
\item Presenter: \textbf{Simon}.
\item Animated build: click 2 times to reveal the remaining steps before moving on.
\item Pause briefly before the bank example — let the concept land first.
\item Close on the takeaway line before advancing.
\end{itemize}
\end{cuebox}
\end{minipage}
\vspace{2mm}
\begin{saybox}
Collect linked evidence and route it to an owner who can investigate and act. \textbf{Signals}: input, execution, and predictive evidence arrive at different times. \textbf{Context}: link observations to run, feature, model, and cohort identity. \textbf{Routing}: an alert has a defined recipient and response policy. \textbf{Authority}: investigation and release approval remain separate. \textit{At our bank,} a CRM missingness alert links to B017 and reaches source / operations owners; it does not promote M2. \textit{A useful monitoring design connects evidence, ownership, and a controlled response.}
\end{saybox}
\par\noindent{\scriptsize\textcolor{teal!60!black}{Reading anchor: Simon’s notes, §2–4 · Bass, Ch. 11 / §14.1.5 · Teaching expansion of monitoring responsibilities}}
\clearpage
\section*{\normalsize 19 / 55 \quad Data drift: a different customer mix}
\noindent\textcolor{teal!70!black}{\small Block 1 / Monitoring \& degradation \textbullet\ Simon}\par\vspace{2mm}
\noindent\begin{minipage}[t]{0.40\textwidth}\vspace{0pt}
\fcolorbox{teal!30}{white}{\includegraphics[width=\dimexpr\linewidth-2\fboxsep-2\fboxrule\relax]{previews/speaker-notes/slide-19.png}}
\end{minipage}\hfill
\begin{minipage}[t]{0.57\textwidth}\vspace{0pt}
\begin{cuebox}
\begin{itemize}
\item Presenter: \textbf{Simon}.
\item Animated build: click 1 time to reveal the remaining steps before moving on.
\item Pause briefly before the bank example — let the concept land first.
\item Close on the takeaway line before advancing.
\end{itemize}
\end{cuebox}
\end{minipage}
\vspace{2mm}
\begin{saybox}
Compare the same input definition across a reference cohort and the current cohort. \textbf{What changes}: the distribution of valid model inputs. \textbf{Bank example}: a campaign attracts more low-activity customers. \textbf{Available now}: compare inputs without waiting for churn labels. \textbf{Limit}: a changed distribution alone does not prove lost accuracy. \textit{At our bank,} both diagrams count 100 fictional customers using the same seven-day login bins. Only the customer mix changes. \textit{Choose a meaningful reference, comparable cohorts, and a defined response.}
\end{saybox}
\par\noindent{\scriptsize\textcolor{teal!60!black}{Reading anchor: Simon’s notes, §2–4 · Bass, Ch. 11 / §14.1.5 · Invented counts}}
\clearpage
\section*{\normalsize 20 / 55 \quad Concept drift: a different relationship}
\noindent\textcolor{teal!70!black}{\small Block 1 / Monitoring \& degradation \textbullet\ Simon}\par\vspace{2mm}
\noindent\begin{minipage}[t]{0.40\textwidth}\vspace{0pt}
\fcolorbox{teal!30}{white}{\includegraphics[width=\dimexpr\linewidth-2\fboxsep-2\fboxrule\relax]{previews/speaker-notes/slide-20.png}}
\end{minipage}\hfill
\begin{minipage}[t]{0.57\textwidth}\vspace{0pt}
\begin{cuebox}
\begin{itemize}
\item Presenter: \textbf{Simon}.
\item Animated build: click 1 time to reveal the remaining steps before moving on.
\item Pause briefly before the bank example — let the concept land first.
\item Close on the takeaway line before advancing.
\end{itemize}
\end{cuebox}
\end{minipage}
\vspace{2mm}
\begin{saybox}
The same inputs can have different outcomes while the deployed model keeps its earlier boundary. \textbf{What changes}: churn risk given the same valid input information. \textbf{Bank example}: new fees make active customers leave more often. \textbf{Evidence}: link predictions to mature, reliable churn outcomes. \textbf{Model status}: M1 stays fixed until a separately evaluated release changes it. \textit{At our bank,} the two panels use the same 15 fictional feature positions. After the hypothetical change, some previously staying customers churn. \textit{A missing field is not concept drift; a statistical alarm is not a retraining approval.}
\end{saybox}
\par\noindent{\scriptsize\textcolor{teal!60!black}{Reading anchor: Simon’s notes, §2–4 · Bass, Ch. 11 / §14.1.5 · Schematic relationships}}
\clearpage
\section*{\normalsize 21 / 55 \quad Concept drift: how M1 starts making errors}
\noindent\textcolor{teal!70!black}{\small Block 1 / Monitoring \& degradation \textbullet\ Simon}\par\vspace{2mm}
\noindent\begin{minipage}[t]{0.40\textwidth}\vspace{0pt}
\fcolorbox{teal!30}{white}{\includegraphics[width=\dimexpr\linewidth-2\fboxsep-2\fboxrule\relax]{previews/speaker-notes/slide-21.png}}
\end{minipage}\hfill
\begin{minipage}[t]{0.57\textwidth}\vspace{0pt}
\begin{cuebox}
\begin{itemize}
\item Presenter: \textbf{Simon}.
\item Animated build: click 1 time to reveal the remaining steps before moving on.
\item Pause briefly before the bank example — let the concept land first.
\item Close on the takeaway line before advancing.
\item Adaptation note on file for this slide; do not read aloud unless asked.
\end{itemize}
\end{cuebox}
\end{minipage}
\vspace{2mm}
\begin{saybox}
A fixed model can miss churners because the world changed—not because its software stopped running. \textbf{Inputs stay fixed}: the two cohorts use identical feature positions. \textbf{M1 stays fixed}: its earlier decision boundary is unchanged. \textbf{Outcomes change}: a different relationship creates missed churners. \textbf{Evidence}: the example uses mature labels, not outcomes already known on scoring day. \textit{At our bank,} c1 has 5 logins and 5 products; C2 has 11 logins and 1 product. The same inputs now accompany different observed churn outcomes. \textit{Changed labels reveal the mismatch. Investigate causes, evaluate candidates, and use controlled release rather than automatic promotion.}
\end{saybox}
\par\noindent{\scriptsize\textcolor{teal!60!black}{Reading anchor: Simon’s notes, §2–4 · Bass, Ch. 11 / §14.1.5 · Invented mature cohorts and deliberately simplified decision rule}}
\clearpage
\section*{\normalsize 22 / 55 \quad Input checks now; quality checks later}
\noindent\textcolor{teal!70!black}{\small Block 1 / Monitoring \& degradation \textbullet\ Simon}\par\vspace{2mm}
\noindent\begin{minipage}[t]{0.40\textwidth}\vspace{0pt}
\fcolorbox{teal!30}{white}{\includegraphics[width=\dimexpr\linewidth-2\fboxsep-2\fboxrule\relax]{previews/speaker-notes/slide-22.png}}
\end{minipage}\hfill
\begin{minipage}[t]{0.57\textwidth}\vspace{0pt}
\begin{cuebox}
\begin{itemize}
\item Presenter: \textbf{Simon}.
\item Animated build: click 2 times to reveal the remaining steps before moving on.
\item Pause briefly before the bank example — let the concept land first.
\item Close on the takeaway line before advancing.
\end{itemize}
\end{cuebox}
\end{minipage}
\vspace{2mm}
\begin{saybox}
The target window creates a delay between a prediction and complete outcome evidence. \textbf{At scoring time}: check freshness, missingness, execution, and coverage. \textbf{During the window}: the full churn outcome is still unknown. \textbf{After ingestion}: join labels to the original score and target window. \textbf{Evaluate by cohort}: precision, recall, and calibration answer different questions. \textit{At our bank,} ana’s D0 score is saved with its context. Complete next-30-day evaluation must wait for that window and ingestion. \textit{An absent label does not mean the customer stayed.}
\end{saybox}
\par\noindent{\scriptsize\textcolor{teal!60!black}{Reading anchor: Simon’s notes, §2–4 · Bass, Ch. 11 / §14.1.5 · Teaching 30-day window}}
\clearpage
\section*{\normalsize 23 / 55 \quad Drift detection: methods and their purpose}
\noindent\textcolor{teal!70!black}{\small Block 1 / Monitoring \& degradation \textbullet\ Simon}\par\vspace{2mm}
\noindent\begin{minipage}[t]{0.40\textwidth}\vspace{0pt}
\fcolorbox{teal!30}{white}{\includegraphics[width=\dimexpr\linewidth-2\fboxsep-2\fboxrule\relax]{previews/speaker-notes/slide-23.png}}
\end{minipage}\hfill
\begin{minipage}[t]{0.57\textwidth}\vspace{0pt}
\begin{cuebox}
\begin{itemize}
\item Presenter: \textbf{Simon}.
\item Pause briefly before the bank example — let the concept land first.
\item Close on the takeaway line before advancing.
\end{itemize}
\end{cuebox}
\end{minipage}
\vspace{2mm}
\begin{saybox}
Methods compare inputs; they do not directly measure business value or model correctness. \textbf{Reference}: name the population, period, and feature definition. \textbf{Thresholds}: calibrate for the application rather than copying a universal rule. \textbf{Segments}: inspect affected cohorts instead of relying only on an aggregate. \textbf{Action}: investigate drift alongside quality, change, and business evidence. \textit{At our bank,} compare the same seven-day login feature before and after a campaign; investigate why the cohort mix changed. \textit{Interpret the signal alongside data validity, sample size, and downstream quality.}
\end{saybox}
\par\noindent{\scriptsize\textcolor{teal!60!black}{Reading anchor: Simon’s notes, §2–4 · Bass, Ch. 11 / §14.1.5 · Simon’s technical atlas}}
\clearpage
\section*{\normalsize 24 / 55 \quad Data-quality faults and training-serving skew}
\noindent\textcolor{teal!70!black}{\small Block 1 / Monitoring \& degradation \textbullet\ Simon}\par\vspace{2mm}
\noindent\begin{minipage}[t]{0.40\textwidth}\vspace{0pt}
\fcolorbox{teal!30}{white}{\includegraphics[width=\dimexpr\linewidth-2\fboxsep-2\fboxrule\relax]{previews/speaker-notes/slide-24.png}}
\end{minipage}\hfill
\begin{minipage}[t]{0.57\textwidth}\vspace{0pt}
\begin{cuebox}
\begin{itemize}
\item Presenter: \textbf{Simon}.
\item Animated build: click 3 times to reveal the remaining steps before moving on.
\item Pause briefly before the bank example — let the concept land first.
\item Close on the takeaway line before advancing.
\end{itemize}
\end{cuebox}
\end{minipage}
\vspace{2mm}
\begin{saybox}
Invalid data and inconsistent preparation differ from a shifted valid population. \textbf{Data-quality fault}: missing, stale, corrupt, or incompatible records. \textbf{Skew}: training and production prepare the feature differently. \textbf{Inspect}: definitions, cutoffs, joins, source changes, and code versions. \textbf{Repair the cause}: retraining on defective inputs can preserve the defect. \textit{At our bank,} a missing engagement field after CRM migration is invalid input. A 30-day production count replacing seven days is skew. \textit{Do not use the word “drift” as a diagnosis for every alarm.}
\end{saybox}
\par\noindent{\scriptsize\textcolor{teal!60!black}{Reading anchor: Simon’s notes, §2–4 · Bass, Ch. 11 / §14.1.5}}
\clearpage
\section*{\normalsize 25 / 55 \quad A publication gate contains a bad batch}
\noindent\textcolor{teal!70!black}{\small Block 1 / Monitoring \& degradation \textbullet\ Simon}\par\vspace{2mm}
\noindent\begin{minipage}[t]{0.40\textwidth}\vspace{0pt}
\fcolorbox{teal!30}{white}{\includegraphics[width=\dimexpr\linewidth-2\fboxsep-2\fboxrule\relax]{previews/speaker-notes/slide-25.png}}
\end{minipage}\hfill
\begin{minipage}[t]{0.57\textwidth}\vspace{0pt}
\begin{cuebox}
\begin{itemize}
\item Presenter: \textbf{Simon}.
\item Animated build: click 2 times to reveal the remaining steps before moving on.
\item Pause briefly before the bank example — let the concept land first.
\item Close on the takeaway line before advancing.
\end{itemize}
\end{cuebox}
\end{minipage}
\vspace{2mm}
\begin{saybox}
The CRM incident illustrates a concrete input fault and a controlled response. \textbf{Detect}: invalid records violate the agreed input contract. \textbf{Contain}: hold the entire B017 batch before publication. \textbf{Investigate}: relate cohort, source change, feature definition, and execution. \textbf{Keep M1}: a source fault does not itself justify a new learned model. \textit{At our bank,} ana’s engagement value is missing, not zero. No new score or completed customer contact is invented. \textit{A gate blocks unsafe use; an alert sends evidence to someone who can investigate.}
\end{saybox}
\par\noindent{\scriptsize\textcolor{teal!60!black}{Reading anchor: Simon’s notes, §2–4 · Bass, Ch. 11 / §14.1.5 · Fictional B017 incident}}
\clearpage
\section*{\normalsize 26 / 55 \quad Respond: investigate, contain, repair, verify}
\noindent\textcolor{teal!70!black}{\small Block 1 / Monitoring \& degradation \textbullet\ Simon}\par\vspace{2mm}
\noindent\begin{minipage}[t]{0.40\textwidth}\vspace{0pt}
\fcolorbox{teal!30}{white}{\includegraphics[width=\dimexpr\linewidth-2\fboxsep-2\fboxrule\relax]{previews/speaker-notes/slide-26.png}}
\end{minipage}\hfill
\begin{minipage}[t]{0.57\textwidth}\vspace{0pt}
\begin{cuebox}
\begin{itemize}
\item Presenter: \textbf{Simon}.
\item Animated build: click 2 times to reveal the remaining steps before moving on.
\item Pause briefly before the bank example — let the concept land first.
\item Close on the takeaway line before advancing.
\end{itemize}
\end{cuebox}
\end{minipage}
\vspace{2mm}
\begin{saybox}
A monitoring signal opens a decision process; it does not select or approve a remedy. \textbf{Diagnose} before choosing repair, rollback, fallback, or training. \textbf{Preserve evidence} and the identity of failed attempts. \textbf{Validate the complete path} after the chosen repair. \textbf{Request a model change} only when evidence supports one. \textit{At our bank,} repair the CRM path, validate a new execution B017-R1, and retain M1 unless a separate model change is justified. \textit{Handoff to Luca: which dependencies made this failure possible?}
\end{saybox}
\par\noindent{\scriptsize\textcolor{teal!60!black}{Reading anchor: Simon’s notes, §2–4 · Bass, Ch. 11 / §14.1.5}}
\clearpage
\section*{\normalsize 27 / 55 \quad Hidden dependencies create technical debt}
\noindent\textcolor{teal!70!black}{\small Block 1 / Design decisions \textbullet\ Luca}\par\vspace{2mm}
\noindent\begin{minipage}[t]{0.40\textwidth}\vspace{0pt}
\fcolorbox{teal!30}{white}{\includegraphics[width=\dimexpr\linewidth-2\fboxsep-2\fboxrule\relax]{previews/speaker-notes/slide-27.png}}
\end{minipage}\hfill
\begin{minipage}[t]{0.57\textwidth}\vspace{0pt}
\begin{cuebox}
\begin{itemize}
\item New block: \textbf{Block 1 / Design decisions}.
\item Presenter: \textbf{Luca}.
\item Animated build: click 3 times to reveal the remaining steps before moving on.
\item Pause briefly before the bank example — let the concept land first.
\item Close on the takeaway line before advancing.
\end{itemize}
\end{cuebox}
\end{minipage}
\vspace{2mm}
\begin{saybox}
An innocent upstream change can alter ML behavior without changing the model. \textbf{Hidden coupling}: features depend on sources owned by other teams. \textbf{Pipeline complexity}: parallel preparation paths can diverge. \textbf{Versioning gaps}: a model artifact alone cannot reproduce a decision. \textbf{Context matters}: architecture decisions depend on requirements and constraints. \textit{At our bank,} cRM changes its engagement export for another purpose; both training and production can inherit the change. \textit{Make data dependencies as reviewable as code dependencies.}
\end{saybox}
\par\noindent{\scriptsize\textcolor{teal!60!black}{Reading anchor: Kalinowski, Ch. 14 · Nazir, §§5–6}}
\clearpage
\section*{\normalsize 28 / 55 \quad Make the feature contract explicit}
\noindent\textcolor{teal!70!black}{\small Block 1 / Design decisions \textbullet\ Luca}\par\vspace{2mm}
\noindent\begin{minipage}[t]{0.40\textwidth}\vspace{0pt}
\fcolorbox{teal!30}{white}{\includegraphics[width=\dimexpr\linewidth-2\fboxsep-2\fboxrule\relax]{previews/speaker-notes/slide-28.png}}
\end{minipage}\hfill
\begin{minipage}[t]{0.57\textwidth}\vspace{0pt}
\begin{cuebox}
\begin{itemize}
\item Presenter: \textbf{Luca}.
\item Pause briefly before the bank example — let the concept land first.
\item Close on the takeaway line before advancing.
\end{itemize}
\end{cuebox}
\end{minipage}
\vspace{2mm}
\begin{saybox}
A contract records meaning and ownership as well as structural schema. \textbf{Semantics}: define what the feature actually measures. \textbf{Time}: specify window, cutoff, and joins consistently. \textbf{Validity}: define types, ranges, freshness, and missingness. \textbf{Ownership}: review upstream changes with affected consumers. \textit{At our bank,} engagement\_7d means app logins during seven days ending at the scoring cutoff; missing must not silently become zero. \textit{Schema checks alone cannot detect every change in meaning.}
\end{saybox}
\par\noindent{\scriptsize\textcolor{teal!60!black}{Reading anchor: Bass, Ch. 14 · Nazir, §§5–6 · Teaching contract}}
\clearpage
\section*{\normalsize 29 / 55 \quad Separate responsibilities where change happens}
\noindent\textcolor{teal!70!black}{\small Block 1 / Design decisions \textbullet\ Luca}\par\vspace{2mm}
\noindent\begin{minipage}[t]{0.40\textwidth}\vspace{0pt}
\fcolorbox{teal!30}{white}{\includegraphics[width=\dimexpr\linewidth-2\fboxsep-2\fboxrule\relax]{previews/speaker-notes/slide-29.png}}
\end{minipage}\hfill
\begin{minipage}[t]{0.57\textwidth}\vspace{0pt}
\begin{cuebox}
\begin{itemize}
\item Presenter: \textbf{Luca}.
\item Animated build: click 2 times to reveal the remaining steps before moving on.
\item Pause briefly before the bank example — let the concept land first.
\item Close on the takeaway line before advancing.
\end{itemize}
\end{cuebox}
\end{minipage}
\vspace{2mm}
\begin{saybox}
Clear interfaces let different concerns evolve without unnecessary relearning. \textbf{Single responsibility}: identify each module’s reason to change. \textbf{Interfaces}: specify data and behavior at each boundary. \textbf{Dependency injection}: provide implementations through explicit contracts. \textbf{Deployment}: logical modules do not automatically require separate services. \textit{At our bank,} a connector fix changes preparation; a contact-frequency rule changes campaign policy; a new M2 changes inference. \textit{Separate concerns first; choose deployment units from the actual trade-offs.}
\end{saybox}
\par\noindent{\scriptsize\textcolor{teal!60!black}{Reading anchor: Kalinowski, Ch. 14 · Nazir, §§5–6}}
\clearpage
\section*{\normalsize 30 / 55 \quad Pipeline pattern: order the responsibilities}
\noindent\textcolor{teal!70!black}{\small Block 1 / Design decisions \textbullet\ Luca}\par\vspace{2mm}
\noindent\begin{minipage}[t]{0.40\textwidth}\vspace{0pt}
\fcolorbox{teal!30}{white}{\includegraphics[width=\dimexpr\linewidth-2\fboxsep-2\fboxrule\relax]{previews/speaker-notes/slide-30.png}}
\end{minipage}\hfill
\begin{minipage}[t]{0.57\textwidth}\vspace{0pt}
\begin{cuebox}
\begin{itemize}
\item Presenter: \textbf{Luca}.
\item Animated build: click 2 times to reveal the remaining steps before moving on.
\item Pause briefly before the bank example — let the concept land first.
\item Close on the takeaway line before advancing.
\end{itemize}
\end{cuebox}
\end{minipage}
\vspace{2mm}
\begin{saybox}
Make preparation and validation steps explicit and testable. \textbf{Explicit stages} clarify transformations and dependencies. \textbf{Stage contracts} make intermediate data understandable. \textbf{Tests} cover local behavior and end-to-end compatibility. \textbf{Failure handling} prevents a partial result from masquerading as a complete batch. \textit{At our bank,} validate CRM records before computing engagement\_7d and invoking M1. \textit{Use structure to expose the dependency, not to hide it in orchestration code.}
\end{saybox}
\par\noindent{\scriptsize\textcolor{teal!60!black}{Reading anchor: Kalinowski, Ch. 14 · Nazir, §§5–6}}
\clearpage
\section*{\normalsize 31 / 55 \quad Strategy pattern: replace one implementation}
\noindent\textcolor{teal!70!black}{\small Block 1 / Design decisions \textbullet\ Luca}\par\vspace{2mm}
\noindent\begin{minipage}[t]{0.40\textwidth}\vspace{0pt}
\fcolorbox{teal!30}{white}{\includegraphics[width=\dimexpr\linewidth-2\fboxsep-2\fboxrule\relax]{previews/speaker-notes/slide-31.png}}
\end{minipage}\hfill
\begin{minipage}[t]{0.57\textwidth}\vspace{0pt}
\begin{cuebox}
\begin{itemize}
\item Presenter: \textbf{Luca}.
\item Animated build: click 3 times to reveal the remaining steps before moving on.
\item Pause briefly before the bank example — let the concept land first.
\item Close on the takeaway line before advancing.
\end{itemize}
\end{cuebox}
\end{minipage}
\vspace{2mm}
\begin{saybox}
A stable interface separates the client from a specific implementation choice. \textbf{Stable interface}: the client relies on a contract, not a specific class. \textbf{Alternative implementations} can be supplied through that interface. \textbf{Compatibility} still includes meaning, not only types. \textbf{Evaluation} and release approval still apply when behavior changes. \textit{At our bank,} replace a feature-preparation implementation behind a stable interface; verify equivalent semantics before scoring. \textit{Interchangeability in code is not equivalence in predictive behavior.}
\end{saybox}
\par\noindent{\scriptsize\textcolor{teal!60!black}{Reading anchor: Kalinowski, Ch. 14 · Nazir, §§5–6}}
\clearpage
\section*{\normalsize 32 / 55 \quad Observer pattern: notify, then investigate}
\noindent\textcolor{teal!70!black}{\small Block 1 / Design decisions \textbullet\ Luca}\par\vspace{2mm}
\noindent\begin{minipage}[t]{0.40\textwidth}\vspace{0pt}
\fcolorbox{teal!30}{white}{\includegraphics[width=\dimexpr\linewidth-2\fboxsep-2\fboxrule\relax]{previews/speaker-notes/slide-32.png}}
\end{minipage}\hfill
\begin{minipage}[t]{0.57\textwidth}\vspace{0pt}
\begin{cuebox}
\begin{itemize}
\item Presenter: \textbf{Luca}.
\item Animated build: click 2 times to reveal the remaining steps before moving on.
\item Pause briefly before the bank example — let the concept land first.
\item Close on the takeaway line before advancing.
\end{itemize}
\end{cuebox}
\end{minipage}
\vspace{2mm}
\begin{saybox}
A change notification decouples the signal producer from interested responders. \textbf{Subject}: emits a meaningful change event. \textbf{Subscribers}: receive notifications without tightly coupling the producer. \textbf{Context}: include cohort, reference, and evidence links. \textbf{Authority}: receiving a notification does not grant release privileges. \textit{At our bank,} a missingness event notifies data and operations owners; it does not instruct the registry to promote M2. \textit{A useful notification requests attention and preserves the decision boundary.}
\end{saybox}
\par\noindent{\scriptsize\textcolor{teal!60!black}{Reading anchor: Kalinowski, Ch. 14 · Nazir, §§5–6}}
\clearpage
\section*{\normalsize 33 / 55 \quad Version the system, not only the model}
\noindent\textcolor{teal!70!black}{\small Block 1 / Design decisions \textbullet\ Luca}\par\vspace{2mm}
\noindent\begin{minipage}[t]{0.40\textwidth}\vspace{0pt}
\fcolorbox{teal!30}{white}{\includegraphics[width=\dimexpr\linewidth-2\fboxsep-2\fboxrule\relax]{previews/speaker-notes/slide-33.png}}
\end{minipage}\hfill
\begin{minipage}[t]{0.57\textwidth}\vspace{0pt}
\begin{cuebox}
\begin{itemize}
\item Presenter: \textbf{Luca}.
\item Animated build: click 2 times to reveal the remaining steps before moving on.
\item Pause briefly before the bank example — let the concept land first.
\item Close on the takeaway line before advancing.
\end{itemize}
\end{cuebox}
\end{minipage}
\vspace{2mm}
\begin{saybox}
A prediction depends on a compatible set of data meaning, code, model, and configuration. \textbf{Data meaning}: record source and feature definitions. \textbf{Code}: identify preparation and scoring revisions. \textbf{Model}: identify artifact, runtime, and configuration. \textbf{Outputs}: connect the compatible set to score and decision records. \textit{At our bank,} m1 alone cannot explain Ana’s score if a CRM migration changed the meaning of its engagement input. \textit{Co-versioning connects system and model lifecycles.}
\end{saybox}
\par\noindent{\scriptsize\textcolor{teal!60!black}{Reading anchor: Kalinowski, Ch. 13–14 · Bass, Ch. 11 / §14.1 · Simon’s notes}}
\clearpage
\section*{\normalsize 34 / 55 \quad Modules or services? Make the trade-off}
\noindent\textcolor{teal!70!black}{\small Block 1 / Design decisions \textbullet\ Luca}\par\vspace{2mm}
\noindent\begin{minipage}[t]{0.40\textwidth}\vspace{0pt}
\fcolorbox{teal!30}{white}{\includegraphics[width=\dimexpr\linewidth-2\fboxsep-2\fboxrule\relax]{previews/speaker-notes/slide-34.png}}
\end{minipage}\hfill
\begin{minipage}[t]{0.57\textwidth}\vspace{0pt}
\begin{cuebox}
\begin{itemize}
\item Presenter: \textbf{Luca}.
\item Pause briefly before the bank example — let the concept land first.
\item Close on the takeaway line before advancing.
\end{itemize}
\end{cuebox}
\end{minipage}
\vspace{2mm}
\begin{saybox}
Independent deployment is useful only when it justifies its operating cost. \textbf{Modules} can provide strong logical boundaries with simpler operation. \textbf{Services} can support independent change and scaling. \textbf{Distribution} adds failure modes and version coordination. \textbf{Evidence} and requirements should drive the choice. \textit{At our bank,} separate preparation from policy logically; deploy them independently only when the bank’s needs justify it. \textit{There is no universal best ML architecture.}
\end{saybox}
\par\noindent{\scriptsize\textcolor{teal!60!black}{Reading anchor: Kalinowski, Ch. 14 · Nazir, §§5–6}}
\clearpage
\section*{\normalsize 35 / 55 \quad Architectural tactics connect concerns to controls}
\noindent\textcolor{teal!70!black}{\small Block 2 / Architectural tactics \textbullet\ TBD}\par\vspace{2mm}
\noindent\begin{minipage}[t]{0.40\textwidth}\vspace{0pt}
\fcolorbox{teal!30}{white}{\includegraphics[width=\dimexpr\linewidth-2\fboxsep-2\fboxrule\relax]{previews/speaker-notes/slide-35.png}}
\end{minipage}\hfill
\begin{minipage}[t]{0.57\textwidth}\vspace{0pt}
\begin{cuebox}
\begin{itemize}
\item New block: \textbf{Block 2 / Architectural tactics}.
\item Presenter: \textbf{TBD}.
\item Pause briefly before the bank example — let the concept land first.
\item Close on the takeaway line before advancing.
\end{itemize}
\end{cuebox}
\end{minipage}
\vspace{2mm}
\begin{saybox}
Use explicit mechanisms to address uncertainty, change, recovery, and auditability. \textbf{Separation} isolates data preparation, model execution, and business policy. \textbf{Observation} makes operational and predictive behavior reviewable. \textbf{Recovery} restores a compatible path rather than only one artifact. \textbf{Governance} links evidence, approval, access, and responsibility. \textit{At our bank,} a validation boundary contains a bad CRM batch; version links and approvals support controlled repair. \textit{A tactic is credible when its mechanism, owner, and evidence are clear.}
\end{saybox}
\par\noindent{\scriptsize\textcolor{teal!60!black}{Reading anchor: Bass, Ch. 14 · Bass, Ch. 9–11 · Group Block 2 outline}}
\clearpage
\section*{\normalsize 36 / 55 \quad CI, CD, and CT need different evidence}
\noindent\textcolor{teal!70!black}{\small Block 2 / DevOps \& MLOps \textbullet\ Matheus}\par\vspace{2mm}
\noindent\begin{minipage}[t]{0.40\textwidth}\vspace{0pt}
\fcolorbox{teal!30}{white}{\includegraphics[width=\dimexpr\linewidth-2\fboxsep-2\fboxrule\relax]{previews/speaker-notes/slide-36.png}}
\end{minipage}\hfill
\begin{minipage}[t]{0.57\textwidth}\vspace{0pt}
\begin{cuebox}
\begin{itemize}
\item New block: \textbf{Block 2 / DevOps \& MLOps}.
\item Presenter: \textbf{Matheus}.
\item Animated build: click 3 times to reveal the remaining steps before moving on.
\item Pause briefly before the bank example — let the concept land first.
\item Close on the takeaway line before advancing.
\end{itemize}
\end{cuebox}
\end{minipage}
\vspace{2mm}
\begin{saybox}
Integration, delivery, deployment, and training automate different responsibilities. \textbf{CI}: integrate changes and run code, contract, and path checks. \textbf{Delivery}: keep a compatible change ready for an approval decision. \textbf{Deployment}: automatically release changes passing the defined policy. \textbf{CT}: train and evaluate candidates; completion is not production approval. \textit{At our bank,} a connector repair can use CI/CD while retaining M1. Retraining is not required for every pipeline change. \textit{Automation executes policy; it does not make acceptance criteria identical.}
\end{saybox}
\par\noindent{\scriptsize\textcolor{teal!60!black}{Reading anchor: Bass, §14.1 · Kalinowski, §16.3–16.5 (supporting)}}
\clearpage
\section*{\normalsize 37 / 55 \quad Candidate training ends before promotion}
\noindent\textcolor{teal!70!black}{\small Block 2 / DevOps \& MLOps \textbullet\ Matheus}\par\vspace{2mm}
\noindent\begin{minipage}[t]{0.40\textwidth}\vspace{0pt}
\fcolorbox{teal!30}{white}{\includegraphics[width=\dimexpr\linewidth-2\fboxsep-2\fboxrule\relax]{previews/speaker-notes/slide-37.png}}
\end{minipage}\hfill
\begin{minipage}[t]{0.57\textwidth}\vspace{0pt}
\begin{cuebox}
\begin{itemize}
\item Presenter: \textbf{Matheus}.
\item Animated build: click 2 times to reveal the remaining steps before moving on.
\item Pause briefly before the bank example — let the concept land first.
\item Close on the takeaway line before advancing.
\end{itemize}
\end{cuebox}
\end{minipage}
\vspace{2mm}
\begin{saybox}
A registry stores candidate evidence; a separate release decision controls production use. \textbf{Reproducibility}: record data, code, configuration, and artifacts. \textbf{Evaluation}: compare candidate evidence against the baseline and relevant cohorts. \textbf{Registration}: preserve identity and lineage for review. \textbf{Promotion}: authorize a specific compatible set with recovery planned. \textit{At our bank,} m2 remains a candidate until reviewed. M1 remains the approved scoring artifact during the CRM repair. \textit{The registry is part of governance, not the whole approval process.}
\end{saybox}
\par\noindent{\scriptsize\textcolor{teal!60!black}{Reading anchor: Bass, §14.1 · Kalinowski, §16.3–16.5 (supporting)}}
\clearpage
\section*{\normalsize 38 / 55 \quad Repair the input path; retain M1 when appropriate}
\noindent\textcolor{teal!70!black}{\small Block 2 / DevOps \& MLOps \textbullet\ Matheus}\par\vspace{2mm}
\noindent\begin{minipage}[t]{0.40\textwidth}\vspace{0pt}
\fcolorbox{teal!30}{white}{\includegraphics[width=\dimexpr\linewidth-2\fboxsep-2\fboxrule\relax]{previews/speaker-notes/slide-38.png}}
\end{minipage}\hfill
\begin{minipage}[t]{0.57\textwidth}\vspace{0pt}
\begin{cuebox}
\begin{itemize}
\item Presenter: \textbf{Matheus}.
\item Animated build: click 2 times to reveal the remaining steps before moving on.
\item Pause briefly before the bank example — let the concept land first.
\item Close on the takeaway line before advancing.
\end{itemize}
\end{cuebox}
\end{minipage}
\vspace{2mm}
\begin{saybox}
The CRM incident requires a verified source repair, not an assumed new model. \textbf{Repair} the source or preparation defect. \textbf{Version} the reviewed fix and compatible configuration. \textbf{Re-run} under a new execution identity and validate the complete path. \textbf{Publish} only after the configured evidence and approval requirements pass. \textit{At our bank,} b017 is the held attempt. B017-R1 links the reviewed fix, compatible M1, validation evidence, and publication approval. \textit{A successful fix needs evidence at the consumer boundary, not only a passing connector test.}
\end{saybox}
\par\noindent{\scriptsize\textcolor{teal!60!black}{Reading anchor: Bass, §14.1 · Kalinowski, §16.3–16.5 (supporting) · Simon’s notes, §4}}
\clearpage
\section*{\normalsize 39 / 55 \quad Trace the prediction and the decision}
\noindent\textcolor{teal!70!black}{\small Block 2 / DevOps \& MLOps \textbullet\ Matheus}\par\vspace{2mm}
\noindent\begin{minipage}[t]{0.40\textwidth}\vspace{0pt}
\fcolorbox{teal!30}{white}{\includegraphics[width=\dimexpr\linewidth-2\fboxsep-2\fboxrule\relax]{previews/speaker-notes/slide-39.png}}
\end{minipage}\hfill
\begin{minipage}[t]{0.57\textwidth}\vspace{0pt}
\begin{cuebox}
\begin{itemize}
\item Presenter: \textbf{Matheus}.
\item Animated build: click 2 times to reveal the remaining steps before moving on.
\item Pause briefly before the bank example — let the concept land first.
\item Close on the takeaway line before advancing.
\end{itemize}
\end{cuebox}
\end{minipage}
\vspace{2mm}
\begin{saybox}
Lineage connects input evidence, execution, approval, and downstream action. \textbf{Prediction identity} links customer, time, and target window. \textbf{Execution evidence} identifies compatible code, features, model, and config. \textbf{Decision evidence} records policy version and action. \textbf{Access controls} protect sensitive reconstruction material. \textit{At our bank,} an investigator can trace Ana’s original score and the campaign’s separate no-contact or contact decision. \textit{A model version alone cannot reconstruct the full decision.}
\end{saybox}
\par\noindent{\scriptsize\textcolor{teal!60!black}{Reading anchor: Bass, §14.1 · Kalinowski, §16.3–16.5 (supporting) · Simon’s notes, §4–6}}
\clearpage
\section*{\normalsize 40 / 55 \quad Recovery restores a compatible path}
\noindent\textcolor{teal!70!black}{\small Block 2 / DevOps \& MLOps \textbullet\ Matheus}\par\vspace{2mm}
\noindent\begin{minipage}[t]{0.40\textwidth}\vspace{0pt}
\fcolorbox{teal!30}{white}{\includegraphics[width=\dimexpr\linewidth-2\fboxsep-2\fboxrule\relax]{previews/speaker-notes/slide-40.png}}
\end{minipage}\hfill
\begin{minipage}[t]{0.57\textwidth}\vspace{0pt}
\begin{cuebox}
\begin{itemize}
\item Presenter: \textbf{Matheus}.
\item Animated build: click 2 times to reveal the remaining steps before moving on.
\item Pause briefly before the bank example — let the concept land first.
\item Close on the takeaway line before advancing.
\end{itemize}
\end{cuebox}
\end{minipage}
\vspace{2mm}
\begin{saybox}
Rehearse restoration before release, and verify it at the consumer boundary. \textbf{Know the target}: identify the last compatible release or validated fallback. \textbf{Rehearse}: test the recovery procedure before relying on it. \textbf{Restore the set}: include data meaning, runtime, and configuration. \textbf{Verify}: inspect stored scores and actual consumer behavior. \textit{At our bank,} rolling M1 back cannot repair a CRM field that is still missing; the input path must be restored. \textit{Handoff to security: who may read, change, approve, or operate these paths?}
\end{saybox}
\par\noindent{\scriptsize\textcolor{teal!60!black}{Reading anchor: Breck, infrastructure tests 6–7 · Simon’s notes, §4}}
\clearpage
\section*{\normalsize 41 / 55 \quad Security threats depend on attacker access}
\noindent\textcolor{teal!70!black}{\small Block 2 / Security \& privacy \textbullet\ Simon + Matheus}\par\vspace{2mm}
\noindent\begin{minipage}[t]{0.40\textwidth}\vspace{0pt}
\fcolorbox{teal!30}{white}{\includegraphics[width=\dimexpr\linewidth-2\fboxsep-2\fboxrule\relax]{previews/speaker-notes/slide-41.png}}
\end{minipage}\hfill
\begin{minipage}[t]{0.57\textwidth}\vspace{0pt}
\begin{cuebox}
\begin{itemize}
\item New block: \textbf{Block 2 / Security \& privacy}.
\item Presenter: \textbf{Simon + Matheus}.
\item Pause briefly before the bank example — let the concept land first.
\item Close on the takeaway line before advancing.
\end{itemize}
\end{cuebox}
\end{minipage}
\vspace{2mm}
\begin{saybox}
Locate the attacker’s control before naming the ML attack mechanism. \textbf{Poisoning}: altered learning inputs influence the trained model. \textbf{Inference manipulation}: altered production inputs affect an existing model’s output. \textbf{Privacy inference}: access can reveal sensitive membership or attributes. \textbf{Assumptions matter}: access, knowledge, and goals differ by threat. \textit{At our bank,} a compromised CRM writer could affect current inputs; those records create a poisoning path only if later used for training. \textit{The accidental CRM migration remains a different scenario.}
\end{saybox}
\par\noindent{\scriptsize\textcolor{teal!60!black}{Reading anchor: Simon’s notes, §5–7 · Bass, Ch. 9–10 / §14.2 · NIST AI 100-2e2025}}
\clearpage
\section*{\normalsize 42 / 55 \quad Enforce access at the resource boundary}
\noindent\textcolor{teal!70!black}{\small Block 2 / Security \& privacy \textbullet\ Simon + Matheus}\par\vspace{2mm}
\noindent\begin{minipage}[t]{0.40\textwidth}\vspace{0pt}
\fcolorbox{teal!30}{white}{\includegraphics[width=\dimexpr\linewidth-2\fboxsep-2\fboxrule\relax]{previews/speaker-notes/slide-42.png}}
\end{minipage}\hfill
\begin{minipage}[t]{0.57\textwidth}\vspace{0pt}
\begin{cuebox}
\begin{itemize}
\item Presenter: \textbf{Simon + Matheus}.
\item Animated build: click 2 times to reveal the remaining steps before moving on.
\item Pause briefly before the bank example — let the concept land first.
\item Close on the takeaway line before advancing.
\end{itemize}
\end{cuebox}
\end{minipage}
\vspace{2mm}
\begin{saybox}
Authentication, authorization, validation, and audit solve different parts of the problem. \textbf{Authenticate}: establish caller identity. \textbf{Authorize}: check operation, resource, and scope. \textbf{Validate}: enforce the input contract independently. \textbf{Audit}: preserve useful access and change evidence. \textit{At our bank,} the scorer may read the approved M1 artifact; it must not gain permission to promote M2 merely by reaching the registry. \textit{A trust boundary needs an enforced policy, not only a box on a diagram.}
\end{saybox}
\par\noindent{\scriptsize\textcolor{teal!60!black}{Reading anchor: Simon’s notes, §5–7 · Bass, Ch. 9–10 / §14.2 · NIST SP 800-207, §3}}
\clearpage
\section*{\normalsize 43 / 55 \quad Least privilege follows responsibilities}
\noindent\textcolor{teal!70!black}{\small Block 2 / Security \& privacy \textbullet\ Simon + Matheus}\par\vspace{2mm}
\noindent\begin{minipage}[t]{0.40\textwidth}\vspace{0pt}
\fcolorbox{teal!30}{white}{\includegraphics[width=\dimexpr\linewidth-2\fboxsep-2\fboxrule\relax]{previews/speaker-notes/slide-43.png}}
\end{minipage}\hfill
\begin{minipage}[t]{0.57\textwidth}\vspace{0pt}
\begin{cuebox}
\begin{itemize}
\item Presenter: \textbf{Simon + Matheus}.
\item Pause briefly before the bank example — let the concept land first.
\item Close on the takeaway line before advancing.
\end{itemize}
\end{cuebox}
\end{minipage}
\vspace{2mm}
\begin{saybox}
Separate source writes, candidate creation, release approval, and campaign use. \textbf{Scope identities} to the operations and resources they need. \textbf{Separate duties}: training does not imply release authority. \textbf{Protect credentials} and avoid secrets in code or broad logs. \textbf{Review changes}: connect access, source updates, and approvals. \textit{At our bank,} a campaign user reads an approved score view, not unrestricted training data or production deployment credentials. \textit{This abbreviated teaching matrix is a starting point, not a complete bank permission design.}
\end{saybox}
\par\noindent{\scriptsize\textcolor{teal!60!black}{Reading anchor: Simon’s notes, §5–7 · Bass, Ch. 9–10 / §14.2 · Proposed permissions}}
\clearpage
\section*{\normalsize 44 / 55 \quad Privacy: expose what each role needs}
\noindent\textcolor{teal!70!black}{\small Block 2 / Security \& privacy \textbullet\ Simon + Matheus}\par\vspace{2mm}
\noindent\begin{minipage}[t]{0.40\textwidth}\vspace{0pt}
\fcolorbox{teal!30}{white}{\includegraphics[width=\dimexpr\linewidth-2\fboxsep-2\fboxrule\relax]{previews/speaker-notes/slide-44.png}}
\end{minipage}\hfill
\begin{minipage}[t]{0.57\textwidth}\vspace{0pt}
\begin{cuebox}
\begin{itemize}
\item Presenter: \textbf{Simon + Matheus}.
\item Animated build: click 3 times to reveal the remaining steps before moving on.
\item Pause briefly before the bank example — let the concept land first.
\item Close on the takeaway line before advancing.
\end{itemize}
\end{cuebox}
\end{minipage}
\vspace{2mm}
\begin{saybox}
Campaign use and incident investigation require different views of customer information. \textbf{Minimize collection} and the information exposed in outputs. \textbf{Separate views}: campaigns and investigators have different purposes. \textbf{Protect logs} and define retention and access rules. \textbf{Preserve accountability} without making raw evidence broadly available. \textit{At our bank,} campaign staff may need Ana’s approved score and eligibility context; investigators may need restricted versioned evidence. \textit{Privacy is an access-and-data-use design, not a label attached to a technique.}
\end{saybox}
\par\noindent{\scriptsize\textcolor{teal!60!black}{Reading anchor: Simon’s notes, §5–7 · Bass, Ch. 9–10 / §14.2}}
\clearpage
\section*{\normalsize 45 / 55 \quad Privacy techniques have specific limits}
\noindent\textcolor{teal!70!black}{\small Block 2 / Security \& privacy \textbullet\ Simon + Matheus}\par\vspace{2mm}
\noindent\begin{minipage}[t]{0.40\textwidth}\vspace{0pt}
\fcolorbox{teal!30}{white}{\includegraphics[width=\dimexpr\linewidth-2\fboxsep-2\fboxrule\relax]{previews/speaker-notes/slide-45.png}}
\end{minipage}\hfill
\begin{minipage}[t]{0.57\textwidth}\vspace{0pt}
\begin{cuebox}
\begin{itemize}
\item Presenter: \textbf{Simon + Matheus}.
\item Pause briefly before the bank example — let the concept land first.
\item Close on the takeaway line before advancing.
\end{itemize}
\end{cuebox}
\end{minipage}
\vspace{2mm}
\begin{saybox}
Choose a mechanism for a defined objective rather than claiming privacy by technique name. \textbf{Specify the objective}: what information should be protected from whom? \textbf{Tokens} and synthetic data do not automatically establish anonymity. \textbf{Federation} changes data location, not every leakage risk. \textbf{DP} needs a defined protected unit and accounting across releases; no derivation here. \textit{At our bank,} protecting one transaction differs from protecting all contributions of a bank customer. \textit{Do not announce a guarantee without its assumptions and evidence.}
\end{saybox}
\par\noindent{\scriptsize\textcolor{teal!60!black}{Reading anchor: Simon’s notes, §5–7 · Bass, Ch. 9–10 / §14.2 · NIST SP 800-226}}
\clearpage
\section*{\normalsize 46 / 55 \quad Use the quality checklist as a review tool}
\noindent\textcolor{teal!70!black}{\small Block 2 / Quality concerns \textbullet\ TBD}\par\vspace{2mm}
\noindent\begin{minipage}[t]{0.40\textwidth}\vspace{0pt}
\fcolorbox{teal!30}{white}{\includegraphics[width=\dimexpr\linewidth-2\fboxsep-2\fboxrule\relax]{previews/speaker-notes/slide-46.png}}
\end{minipage}\hfill
\begin{minipage}[t]{0.57\textwidth}\vspace{0pt}
\begin{cuebox}
\begin{itemize}
\item New block: \textbf{Block 2 / Quality concerns}.
\item Presenter: \textbf{TBD}.
\item Pause briefly before the bank example — let the concept land first.
\item Close on the takeaway line before advancing.
\end{itemize}
\end{cuebox}
\end{minipage}
\vspace{2mm}
\begin{saybox}
Connect concerns to architectural controls, ownership, and observable evidence. \textbf{System and process concerns} both matter in ML-intensive systems. \textbf{Select relevant checks} rather than treating a short list as the entire checklist. \textbf{Map each concern} to a control, owner, and evidence. \textbf{Keep limits visible}: the separate corrigendum remains unread. \textit{At our bank,} completeness validation addresses input quality; it is not by itself evidence of SS5 cyberattack defense. \textit{A checklist supports structured review; it is not a certification claim.}
\end{saybox}
\par\noindent{\scriptsize\textcolor{teal!60!black}{Reading anchor: Bucaioni, published Table 1, pp. 6–7 · Bass, Ch. 14 · Corrigendum contents not reviewed}}
\clearpage
\section*{\normalsize 47 / 55 \quad Predictive quality and system performance}
\noindent\textcolor{teal!70!black}{\small Block 2 / Quality concerns \textbullet\ TBD}\par\vspace{2mm}
\noindent\begin{minipage}[t]{0.40\textwidth}\vspace{0pt}
\fcolorbox{teal!30}{white}{\includegraphics[width=\dimexpr\linewidth-2\fboxsep-2\fboxrule\relax]{previews/speaker-notes/slide-47.png}}
\end{minipage}\hfill
\begin{minipage}[t]{0.57\textwidth}\vspace{0pt}
\begin{cuebox}
\begin{itemize}
\item Presenter: \textbf{TBD}.
\item Pause briefly before the bank example — let the concept land first.
\item Close on the takeaway line before advancing.
\end{itemize}
\end{cuebox}
\end{minipage}
\vspace{2mm}
\begin{saybox}
Choose an acceptable model and system from several quality objectives. \textbf{Accuracy} alone is an incomplete objective, especially with imbalanced outcomes. \textbf{Performance} includes the end-to-end deadline and resource use. \textbf{Complexity} must justify training and operating costs. \textbf{Acceptance} depends on the task and the consequences of errors. \textit{At our bank,} a heavier churn model is useful only if its quality benefit justifies cost and still meets the campaign’s deadline. \textit{Define acceptance criteria before selecting a more complex model.}
\end{saybox}
\par\noindent{\scriptsize\textcolor{teal!60!black}{Reading anchor: Bucaioni, published Table 1, pp. 6–7 · Bass, Ch. 14 · Kalinowski, §13.4}}
\clearpage
\section*{\normalsize 48 / 55 \quad Robustness: invalid, unfamiliar, or ordinary?}
\noindent\textcolor{teal!70!black}{\small Block 2 / Quality concerns \textbullet\ TBD}\par\vspace{2mm}
\noindent\begin{minipage}[t]{0.40\textwidth}\vspace{0pt}
\fcolorbox{teal!30}{white}{\includegraphics[width=\dimexpr\linewidth-2\fboxsep-2\fboxrule\relax]{previews/speaker-notes/slide-48.png}}
\end{minipage}\hfill
\begin{minipage}[t]{0.57\textwidth}\vspace{0pt}
\begin{cuebox}
\begin{itemize}
\item Presenter: \textbf{TBD}.
\item Pause briefly before the bank example — let the concept land first.
\item Close on the takeaway line before advancing.
\end{itemize}
\end{cuebox}
\end{minipage}
\vspace{2mm}
\begin{saybox}
Different inputs need different responses; do not treat every unusual value as the same failure. \textbf{Contract validity} and model applicability are distinct. \textbf{OOD inputs} may be structurally valid but poorly represented in training. \textbf{Confidence} alone is not a guaranteed detector of unfamiliar inputs. \textbf{Fallbacks} must be validated for the intended situation. \textit{At our bank,} a missing activity value is invalid; a new customer population may be valid but outside the model’s reliable coverage. \textit{An architecture needs an explicit response when the model should not be trusted.}
\end{saybox}
\par\noindent{\scriptsize\textcolor{teal!60!black}{Reading anchor: Bucaioni, published Table 1, pp. 6–7 · Bass, Ch. 14 · Simon’s notes, §2–4}}
\clearpage
\section*{\normalsize 49 / 55 \quad Explainability and auditability differ}
\noindent\textcolor{teal!70!black}{\small Block 2 / Quality concerns \textbullet\ TBD}\par\vspace{2mm}
\noindent\begin{minipage}[t]{0.40\textwidth}\vspace{0pt}
\fcolorbox{teal!30}{white}{\includegraphics[width=\dimexpr\linewidth-2\fboxsep-2\fboxrule\relax]{previews/speaker-notes/slide-49.png}}
\end{minipage}\hfill
\begin{minipage}[t]{0.57\textwidth}\vspace{0pt}
\begin{cuebox}
\begin{itemize}
\item Presenter: \textbf{TBD}.
\item Animated build: click 2 times to reveal the remaining steps before moving on.
\item Pause briefly before the bank example — let the concept land first.
\item Close on the takeaway line before advancing.
\end{itemize}
\end{cuebox}
\end{minipage}
\vspace{2mm}
\begin{saybox}
Explaining feature influence and reconstructing a historical decision are complementary tasks. \textbf{Explanation} helps interpret model behavior in context. \textbf{Auditability} needs historical evidence and traceable versions. \textbf{Decision review} includes business policy and action, not only the model. \textbf{Limits}: SHAP and LIME do not prove causality or fairness. \textit{At our bank,} an engagement contribution may help explain Ana’s score; it cannot prove why she would leave the bank. \textit{Protect evidence and test whether authorized reviewers can actually retrieve it.}
\end{saybox}
\par\noindent{\scriptsize\textcolor{teal!60!black}{Reading anchor: Bass, Ch. 14 · Simon’s notes, §6–7}}
\clearpage
\section*{\normalsize 50 / 55 \quad Fairness: prediction versus treatment}
\noindent\textcolor{teal!70!black}{\small Block 2 / Quality concerns \textbullet\ TBD}\par\vspace{2mm}
\noindent\begin{minipage}[t]{0.40\textwidth}\vspace{0pt}
\fcolorbox{teal!30}{white}{\includegraphics[width=\dimexpr\linewidth-2\fboxsep-2\fboxrule\relax]{previews/speaker-notes/slide-50.png}}
\end{minipage}\hfill
\begin{minipage}[t]{0.57\textwidth}\vspace{0pt}
\begin{cuebox}
\begin{itemize}
\item Presenter: \textbf{TBD}.
\item Animated build: click 3 times to reveal the remaining steps before moving on.
\item Pause briefly before the bank example — let the concept land first.
\item Close on the takeaway line before advancing.
\end{itemize}
\end{cuebox}
\end{minipage}
\vspace{2mm}
\begin{saybox}
Name what is being compared before selecting a metric or declaring an improvement. \textbf{Prediction fairness}: compare errors or calibration across relevant groups. \textbf{Treatment fairness}: inspect access to offers and policy effects. \textbf{Context}: select targets based on harms and requirements. \textbf{Trade-offs}: one metric does not settle every fairness concern. \textit{At our bank,} equal churn prediction errors do not imply equal access to retention offers under campaign allocation rules. \textit{Keep predictive quality and the consequences of intervention distinct.}
\end{saybox}
\par\noindent{\scriptsize\textcolor{teal!60!black}{Reading anchor: Bucaioni, published Table 1, pp. 6–7 · Bass, Ch. 14 · Bass, Ch. 10 / §14.2}}
\clearpage
\section*{\normalsize 51 / 55 \quad Airflow: make workflow dependencies visible}
\noindent\textcolor{teal!70!black}{\small Block 2 / Practical tools \textbullet\ TBD}\par\vspace{2mm}
\noindent\begin{minipage}[t]{0.40\textwidth}\vspace{0pt}
\fcolorbox{teal!30}{white}{\includegraphics[width=\dimexpr\linewidth-2\fboxsep-2\fboxrule\relax]{previews/speaker-notes/slide-51.png}}
\end{minipage}\hfill
\begin{minipage}[t]{0.57\textwidth}\vspace{0pt}
\begin{cuebox}
\begin{itemize}
\item New block: \textbf{Block 2 / Practical tools}.
\item Presenter: \textbf{TBD}.
\item Animated build: click 2 times to reveal the remaining steps before moving on.
\item Pause briefly before the bank example — let the concept land first.
\item Close on the takeaway line before advancing.
\end{itemize}
\end{cuebox}
\end{minipage}
\vspace{2mm}
\begin{saybox}
Orchestration supports task order, failure handling, and execution evidence. \textbf{Architecture role}: order dependent pipeline work. \textbf{Evidence}: task state, timing, and run context. \textbf{Action}: operators and source owners inspect failed dependencies. \textbf{Limit}: orchestration does not supply predictive or release approval evidence by itself. \textit{At our bank,} the required CRM validation must sit upstream of scoring and export so its failure blocks dependent work. \textit{Connect each tool to a responsibility, evidence, and a response owner.}
\end{saybox}
\par\noindent{\scriptsize\textcolor{teal!60!black}{Reading anchor: Bass, Ch. 14 · Illustrative Airflow role}}
\clearpage
\section*{\normalsize 52 / 55 \quad MLflow: connect experiments and artifacts}
\noindent\textcolor{teal!70!black}{\small Block 2 / Practical tools \textbullet\ TBD}\par\vspace{2mm}
\noindent\begin{minipage}[t]{0.40\textwidth}\vspace{0pt}
\fcolorbox{teal!30}{white}{\includegraphics[width=\dimexpr\linewidth-2\fboxsep-2\fboxrule\relax]{previews/speaker-notes/slide-52.png}}
\end{minipage}\hfill
\begin{minipage}[t]{0.57\textwidth}\vspace{0pt}
\begin{cuebox}
\begin{itemize}
\item Presenter: \textbf{TBD}.
\item Animated build: click 2 times to reveal the remaining steps before moving on.
\item Pause briefly before the bank example — let the concept land first.
\item Close on the takeaway line before advancing.
\end{itemize}
\end{cuebox}
\end{minipage}
\vspace{2mm}
\begin{saybox}
Model tracking supports reproducibility and review when evidence links are designed explicitly. \textbf{Architecture role}: connect training evidence and model artifacts. \textbf{Evidence}: runs, configuration, metrics, and artifact identity. \textbf{Action}: reviewers compare candidates and verify compatible releases. \textbf{Limit}: no single tracking tool provides every governance control. \textit{At our bank,} link the approved M1 to its evidence while recording B017-R1’s separate input and execution context. \textit{Tool records support the release decision; they do not replace its policy.}
\end{saybox}
\par\noindent{\scriptsize\textcolor{teal!60!black}{Reading anchor: Bass, Ch. 14 · Illustrative MLflow role}}
\clearpage
\section*{\normalsize 53 / 55 \quad Validation and drift reports support diagnosis}
\noindent\textcolor{teal!70!black}{\small Block 2 / Practical tools \textbullet\ TBD}\par\vspace{2mm}
\noindent\begin{minipage}[t]{0.40\textwidth}\vspace{0pt}
\fcolorbox{teal!30}{white}{\includegraphics[width=\dimexpr\linewidth-2\fboxsep-2\fboxrule\relax]{previews/speaker-notes/slide-53.png}}
\end{minipage}\hfill
\begin{minipage}[t]{0.57\textwidth}\vspace{0pt}
\begin{cuebox}
\begin{itemize}
\item Presenter: \textbf{TBD}.
\item Animated build: click 2 times to reveal the remaining steps before moving on.
\item Pause briefly before the bank example — let the concept land first.
\item Close on the takeaway line before advancing.
\end{itemize}
\end{cuebox}
\end{minipage}
\vspace{2mm}
\begin{saybox}
A report is useful when its reference, feature meaning, cohort, and response are clear. \textbf{Validation}: test declared expectations before use. \textbf{Drift comparison}: inspect change against a named reference. \textbf{Tool choice}: reports can support these roles; keep policy outside the screenshot. \textbf{Action}: preserve evidence and assign a response owner. \textit{At our bank,} a missing CRM feature needs a contract-failure response; a changed valid population needs a different investigation. \textit{Tools make evidence easier to collect; architecture defines what that evidence means.}
\end{saybox}
\par\noindent{\scriptsize\textcolor{teal!60!black}{Reading anchor: Simon’s notes, §2–4 · Bass, Ch. 11 / §14.1.5 · Group tool-showcase topic}}
\clearpage
\section*{\normalsize 54 / 55 \quad ML architecture connects five responsibilities}
\noindent\textcolor{teal!70!black}{\small Block 2 / Practical tools \textbullet\ All presenters}\par\vspace{2mm}
\noindent\begin{minipage}[t]{0.40\textwidth}\vspace{0pt}
\fcolorbox{teal!30}{white}{\includegraphics[width=\dimexpr\linewidth-2\fboxsep-2\fboxrule\relax]{previews/speaker-notes/slide-54.png}}
\end{minipage}\hfill
\begin{minipage}[t]{0.57\textwidth}\vspace{0pt}
\begin{cuebox}
\begin{itemize}
\item Presenter: \textbf{All presenters}.
\item Animated build: click 4 times to reveal the remaining steps before moving on.
\item Pause briefly before the bank example — let the concept land first.
\item Close on the takeaway line before advancing.
\end{itemize}
\end{cuebox}
\end{minipage}
\vspace{2mm}
\begin{saybox}
Return to the full system and the choices that make a prediction usable. \textbf{Data}: preserve meaning, validity, and temporal consistency. \textbf{Models}: evaluate uncertainty and manage compatible approved artifacts. \textbf{Operations}: observe, diagnose, release, and recover. \textbf{Decisions}: protect access and govern appropriate business action. \textit{At our bank,} ana’s score becomes usable through validated inputs, traceable execution, an approved release, and a separate campaign policy. \textit{Day 2: foundation models and agents retain these control needs, while components and failure modes change.}
\end{saybox}
\par\noindent{\scriptsize\textcolor{teal!60!black}{Reading anchor: Bass, Ch. 14 · Group synthesis}}
\clearpage
\section*{\normalsize 55 / 55 \quad Sources and adaptation boundaries}
\noindent\textcolor{teal!70!black}{\small Block 2 / Practical tools \textbullet\ All presenters}\par\vspace{2mm}
\noindent\begin{minipage}[t]{0.40\textwidth}\vspace{0pt}
\fcolorbox{teal!30}{white}{\includegraphics[width=\dimexpr\linewidth-2\fboxsep-2\fboxrule\relax]{previews/speaker-notes/slide-55.png}}
\end{minipage}\hfill
\begin{minipage}[t]{0.57\textwidth}\vspace{0pt}
\begin{cuebox}
\begin{itemize}
\item Presenter: \textbf{All presenters}.
\item Pause briefly before the bank example — let the concept land first.
\item Close on the takeaway line before advancing.
\end{itemize}
\end{cuebox}
\end{minipage}
\vspace{2mm}
\begin{saybox}
The bibliography supplies concepts; the bank diagrams and incidents are explicit teaching applications. \textbf{Core readings}: architectural decisions, serving, maintenance, and quality concerns. \textbf{Supporting readings} extend monitoring, governance, and security. \textbf{Illustrations}: Ana, scores, dates, batches, and incidents are invented. \textbf{Source limit}: the separately catalogued checklist corrigendum remains unread. \textit{At our bank,} no new churn outcome, repaired score, retention benefit, or real bank tool execution is claimed. \textit{The separate guide records reading anchors and adaptation limits.}
\end{saybox}
\par\noindent{\scriptsize\textcolor{teal!60!black}{Reading anchor: See source-review.md and presenter guide}}
\clearpage
\end{document}